\section{Simpler equations}
It makes sense to first start simulations with simpler equations where we take $\langle \dot{\theta} \rangle = \langle (\nabla \theta)^2 \rangle = 0$ and forget about the third equation entirely.

\begin{equation}
    \ddot{\phi} + 3H \dot{\phi}
    - \frac{\kappa V_0}{b M_{\rm Pl}}
    \exp\!\left[\frac{\phi}{b M_{\rm Pl}} - \kappa e^{\frac{\phi}{b M_{\rm Pl}}}\right]
    + g^2 \langle \chi^2 \rangle (\phi - \phi_{\rm SBP}) = 0
\end{equation}

\begin{equation}
    \ddot{\tilde{\chi}}_{\mathbf{k}} + 3H \dot{\tilde{\chi}}_{\mathbf{k}}
    + \left[
    \frac{k^2}{a^2}
    + g^2 (\phi - \phi_{\rm SBP})^2
    - \lambda f^2
    + 3 \lambda \langle \chi^2 \rangle
    \right]
    \tilde{\chi}_{\mathbf{k}} = 0
\end{equation}

where we can go a step further and first consider the homogeneous equation

\begin{equation}
    \ddot{\phi} + 3H \dot{\phi}
    - \frac{\kappa V_0}{b M_{\rm Pl}}
    \exp\!\left[\frac{\phi}{b M_{\rm Pl}} - \kappa e^{\frac{\phi}{b M_{\rm Pl}}}\right] = 0
\end{equation}

To compute the averages we must discretise the integral over $\textbf{k}$ modes, and first we will not consider any 2PI corrections

\begin{equation*}
    \langle \chi^2 \rangle = \int_{\mathbb{R}^3} \frac{d^3 p }{(2 \pi)^3} |\tilde{\chi}_{\textbf{p}} (t)|^2 \approx \sum_{\textbf{p}} \omega_\textbf{p} |\tilde{\chi}_\textbf{p} (t)|^2 
\end{equation*}

where $\omega_\textbf{p}$ being the weights. There are several choices of $\omega_{\textbf{p}}$, but since $\langle \chi^2 \rangle$ is not isometric we use
\begin{equation*}
    \omega_{\textbf{k}} = \frac{\delta k_x  \delta k_y \delta k_z}{ (2 \pi)^3 }
\end{equation*}

 which is determined at the start of the simulation when we choose the change of $k$, assuming that the range is divided in equal steps that is $\delta k_x = \delta k_y = \delta k_z = \delta k \equiv \text{const}$. 
\begin{equation*}
    \omega_{\mathbf{k}} = \left( \frac{\delta k}{2 \pi} \right)^3 = \omega
\end{equation*}
yielding the simplefied expression
\begin{equation*}
    \langle \chi^2 \rangle \approx \omega \sum_{\textbf{p}} |\tilde{\chi}_\textbf{p} (t)|^2
\end{equation*}

Another important thing to note is that the hubble parameter $H$ also changes, since it depends on the energy density $\rho$ of the fields
\begin{align*}
    \rho_{\phi} (t) &= \frac{1}{2} \dot{ \phi }^2 + V( \phi ) \\
    \rho_{\chi} (t) &= \int \frac{d^3 k}{(2 \pi)^3} \frac{1}{2} (|\dot{\tilde{\chi}}_{\textbf{k}}|^2 + \omega_{k})
\end{align*}

Then $\rho_{\text{tot}} = \rho_{\phi} + \rho_{\chi}$, however at first we can neglect the $\rho_{\chi}$ term and the hubble parameter thus becomes

\begin{equation*}
    H \approx \sqrt{\frac{1} {3 M_{\text{pl}}^2} \left( \frac{1}{2} \dot{ \phi }^2 + V( \phi ) \right)} = \sqrt{\frac{1} {3 M_{\text{pl}}^2} \left( \frac{1}{2} \dot{ \phi }^2 + V_0 \exp \left[ - \kappa e^{\frac{\phi}{b M_{\text{Pl}}}} \right] \right)}
\end{equation*}

Note that we are using the potential $V(\phi) = V_0 \exp \left[ - \kappa e^{\frac{\phi}{b M_{\text{Pl}}}} \right]$.

Next, I need to determine initial conditions. It makes sense to integrate from the symetry breaking point $\phi_{\rm SBP}$. During kination (point where energy density is dominated by kinetic energy) the $\phi$ field equation simplifies emensly:

\begin{equation*}
    \ddot{\phi} + 3 H \dot{\phi} = 0
\end{equation*}

with a simplified Hubble parameter $ H \approx \frac{|\dot{\phi}|}{\sqrt{6} M_{\rm pl}} $. We can then solve these equations. Note that we will take $t_i = t_{SBP} = 0$ and $v=\dot{\phi}_{\rm SBP}$

\begin{align*}
    \phi = \sqrt{ \frac{2}{3}} \text{sign}(v) \ln \left( 1+\sqrt{\frac{3}{2}\frac{|v| t} {M_{\rm pl}}} \right) M_{\rm pl} \\
    \dot{ \phi } = v \exp( - \sqrt{\frac{3}{2}} \frac{\phi}{M_{\rm pl} \text{sign}(v)})
\end{align*}

That means that initial conditions $\phi_0 = \phi(t_i) = 0$ and $\dot{\phi}_{0} = \dot{\phi}(0) = v$, where I think I can choose $v$ - being a free parameter. We can similarly consider modes for $\chi$ field. For lange $|\phi|$ they satisfy

\begin{equation*}
    \ddot{\chi}_k + 3 H \dot{\chi}_k + \omega^2_k \chi_k = 0
\end{equation*}

where $ \omega^2 = k^2 + g^2 \phi^2 - \lambda f^2 $. The initial conditions are given by the Bunch-Davies vacuum state

\begin{equation*}
    \chi_{\phi, {\rm vac}} = \frac{a^{-1}}{\sqrt{2 \omega_k}} \exp{-i \int^t \omega_k dt^{\prime}}
\end{equation*}

But again, since our initial conditions have $t=0$ the above simplifies to just:

\begin{equation*}
    \chi_k(t_i) = \frac{1}{a(t_i) \sqrt{ \omega_k}} = \frac{1}{\sqrt{ \omega_k }}
\end{equation*}

Where we have used $a(t_i) = 1$. We can also differentiate $\chi_k$ modes to obtain initial conditions for $\dot{\chi}_k$, where we will only keep the adiabatic terms:

\begin{equation*}
    \dot{\chi} = [-H(t_i) - i \omega_k(t_i)] \chi_k(t_i) \approx 0
\end{equation*}

All of these initial conditions are dependent on parameters $g$, $\lambda$, $f$, $v$ and technically $t_i$, though we set it to 0 previously. We can discuss the ranges of these parameters:

\begin{enumerate}
    \item $10^{-5} < \lambda < 1$
    \item $f = 5 \times 10^-7$
    \item $10^{-14} < g < 1$
    \item v is free
\end{enumerate}